% TOP_PROBLEM: {"id": 2305044, "title": "Research Problems in Function Theory — Problem 5.44", "category_id": 8, "category": {"id": 8, "name": "analysis", "display_name": "Analysis", "description": "Limits, continuity, calculus, and function theory.", "slug": "analysis", "order_index": 8, "created_at": "2026-07-31T15:26:25.671Z"}, "status": "open"}
% FIRST_POSTED: null
% ENTRY_KIND: statement_only
% Imported from: batch14.tex; UnsolvedMath ID 2305044
% Compile twice: lualatex 2305044.tex
\documentclass[11pt]{article}
\usepackage{fontspec}
\setmainfont{Latin Modern Roman}
\usepackage{amsmath,amssymb,mathtools,unicode-math}
\setmathfont{Latin Modern Math}
\usepackage{xurl,hyperref}
\hypersetup{hidelinks,pdftitle={UnsolvedMath Problem 2305044}}
\newcommand{\sref}[2]{\hyperlink{source:#1:#2}{[S#2]}}
\newcommand{\cataloguescope}[1]{\paragraph{Mathematical question.}#1}
\begin{document}
% UnsolvedMath ID 2305044; source catalogue code AMR-022-5044
\setcounter{section}{2305043}
\section{Research Problems in Function Theory — Problem 5.44}\label{2305044}
{\small\sffamily UnsolvedMath ID 2305044 · Analysis}
\subsection{Definitions and mathematical statement}

\paragraph{Shared notation.}$\mathbb N=\{1,2,\ldots\}$, and $\mathbb Z,\mathbb Q,\mathbb R,\mathbb C$ denote the integers, rationals, reals and complex numbers. Any other conventions are specified locally.


Write $\mathbb D=\{z\in\mathbb C:|z|<1\}$ and $\mathbb T=\partial\mathbb D$. All Taylor coefficients below are taken at $0$.For $\alpha,\beta>0$ and $|x|=1$, define coefficients by the analytic branches normalized to equal $1$ at $z=0$:
\[(1+xz)^\alpha(1-z)^{-\beta}=\sum_{m\ge0}A_m^{\alpha,\beta}(x)z^m.
\]
Equivalently, $A_m^{\alpha,\beta}(x)=\sum_{j=0}^m\binom\alpha j x^j(\beta)_{m-j}/(m-j)!$, where $\binom\alpha j=\alpha(\alpha-1)\cdots(\alpha-j+1)/j!$, $(\beta)_k=\beta(\beta+1)\cdots(\beta+k-1)$, and both empty products are $1$. For $\alpha=3/2$ and $\beta=1/100$ one has $A_3^{\alpha,\beta}(1)=-95583/2000000<0$, so a bound by this signed number already fails at $x=1$.
\paragraph{Statement recorded in the supplied source.}
For $0<\alpha<1$ and $\beta=1$, does
\[|A_{2n+1}^{\alpha,1}(x)|\le |A_{2n+1}^{\alpha,1}(1)|\qquad(n\ge2,\ |x|=1)\]
always hold? More generally, investigate the same odd-coefficient domination for arbitrary $\alpha,\beta>0$:
\[|A_{2n+1}^{\alpha,\beta}(x)|\le |A_{2n+1}^{\alpha,\beta}(1)|\qquad(n\ge2,\ |x|=1).\]
The source separately poses the literal first general odd-coefficient inequality
\[|A_3^{\alpha,\beta}(x)|\le A_3^{\alpha,\beta}(1)\qquad(\alpha,\beta>0,\ |x|=1),\]
with a signed right-hand side; this historical subsidiary inequality is distinct from the modulus comparison above. \sref{2305044}{2}
\subsection{Short English statement}
Does the coefficient at $x=1$ maximize the modulus of the relevant odd coefficients, first for the original $\alpha$ range and then for arbitrary positive $\alpha$ and $\beta$? The separately posed cubic-coefficient inequality uses a signed right-hand side and fails for some positive parameters.
\subsection{Sources}
\begin{description}
\item[\hypertarget{source:2305044:1}{S1}] ULAM.AI and the UnsolvedMath Contributors, UnsolvedMath: A Curated Collection of Open Mathematics Problems, record 2305044 (AMR-022-5044). CC BY 4.0. \url{https://huggingface.co/datasets/ulamai/UnsolvedMath}.
\item[\hypertarget{source:2305044:2}{S2}] W. K. Hayman and E. F. Lingham, Research Problems in Function Theory (2018), arXiv:1809.07200, Chapter 5, Problem 5.44, its update and the preceding shared notation. \url{https://arxiv.org/abs/1809.07200}.
\end{description}
\label{end:2305044}
\end{document}
